\chapter{Цілеспрямовані сталі загрози та полювання за загрозами}
\chaptershort{APT та полювання за загрозами}
\label{chap:ch13}

Масові автоматизовані загрози (розділ~\ref{chap:ch01}) протистоять автоматизованим засобам захисту (розділи~\ref{chap:ch06}--\ref{chap:ch11}). Проте найнебезпечніший супротивник --- \term{цілеспрямована стала загроза} (APT, advanced persistent threat): добре ресурсоване угруповання, що наполегливо, тривало й адаптивно досягає конкретних цілей попри захист. Проти APT реактивне виявлення недостатнє: такий супротивник уникає відомих сигнатур, зловживає легітимними інструментами (LotL, розділ~\ref{chap:ch11}) і може роками лишатися непоміченим. Відповіддю є \term{полювання за загрозами} (threat hunting) --- проактивний, керований гіпотезами пошук прихованих загроз, що \emph{не} спрацювали на автоматичному виявленні. Розділ розглядає феномен APT і його життєвий цикл, поняття часу присутності (dwell time), парадигму полювання за загрозами, його методологію й гіпотези, бокове переміщення й living-off-the-land як об'єкти полювання, зрілість і метрики, і завершується кейсом полювання за прихованим угрупованням. Полювання спирається на телеметрію (розділ~\ref{chap:ch10}), розвідку загроз (розділ~\ref{chap:ch12}) і живить реагування (розділ~\ref{chap:ch16}).

\section{Феномен APT та його життєвий цикл}\label{sec:ch13-apt}

\term{APT} характеризують три ознаки: \emph{advanced} (високі можливості --- власні інструменти, експлойти, ресурси); \emph{persistent} (наполегливість --- угруповання не відступає після невдачі, а адаптується й повертається); \emph{threat} (цілеспрямованість --- за атакою стоїть організація з конкретними цілями, часто державна або пов'язана з державою). Перші задокументовані кейси кіберспіонажу~\cite{ch13:stoll} та відкриття діяльності державних угруповань~\cite{ch13:mandiant} показали, що такі актори діють методично, за життєвим циклом ланцюга ураження (розділ~\ref{chap:ch04}): розвідка цілі, початковий доступ, закріплення, підвищення привілеїв, бокове переміщення, досягнення мети (шпигунство, саботаж), приховування слідів.

Відмінність APT від масового ЗПК --- не стільки технічна витонченість кожного кроку, скільки \emph{цілеспрямованість і наполегливість}. Угруповання вивчає конкретну жертву, адаптує інструменти під її середовище, діє повільно й обережно, зливаючись із легітимною активністю, щоб уникнути виявлення. Саме тому ранні аналізи~\cite{ch13:virvilis} наголошували: типовий периметровий захист і сигнатурне виявлення систематично пропускають APT, бо ті не застосовують відомого шкідливого коду там, де його чекають. Відмінності APT від масових загроз узагальнено в табл.~\thechapter.2.

\tabhead{Таблиця~\thechapter.2 --- APT та масовий ЗПК: порівняння}

\begin{center}\footnotesize
\begin{tabularx}{\linewidth}{L{20mm} Y Y}
\toprule
Ознака & Масовий ЗПК & APT \\
\midrule
Ціль & будь-хто вразливий & конкретна жертва \\
Мотив & миттєва вигода & стратегічна мета \\
Наполегливість & разова спроба & адаптація, повернення \\
Інструменти & типові, відомі & власні, LotL \\
Виявлення & сигнатури ловлять & потребує полювання \\
\bottomrule
\end{tabularx}
\end{center}

Варто застерегти від спрощеного ототожнення APT з надскладними технологіями. Угруповання застосовує рівно стільки витонченості, скільки потрібно: якщо ціль беруть фішингом і викраденими паролями, дорогі експлойти нульового дня притримують. Саме економність і прагматизм роблять APT небезпечним --- він не «світить» рідкісними інструментами там, де вистачає буденних, а отже й не лишає помітних слідів. Це пояснює, чому наголос у протидії зміщується з пошуку екзотичного шкідливого коду на виявлення \emph{аномальної поведінки} цілком звичайних засобів.

\section{Час присутності (dwell time)}\label{sec:ch13-dwell}

Ключовий показник ефективності захисту проти APT --- \term{час присутності} (dwell time): проміжок від компрометації до виявлення. Велике dwell time означає, що супротивник довго діяв непоміченим, встигнувши досягти цілей. Складові часу реагування формалізовано у~\eqref{eq:ch10-dwell}; стратегічна мета захисту --- мінімізувати dwell time, бо збитки від APT зростають із часом присутності.

\begin{tikzfig}
\begin{tikzpicture}[font=\scriptsize,>={Stealth[length=1.5mm]}]
\draw[->] (0,0) -- (8,0) node[right,font=\tiny]{час};
\foreach \x/\l in {0.4/компрометація, 2.2/закріплення, 4.0/бокове перем., 5.8/дія на цілі} {
  \draw (\x,0.1)--(\x,-0.1); \node[font=\tiny,rotate=25,anchor=west] at (\x,0.2) {\l};}
\draw[<->,thick] (0.4,-0.6)--(6.5,-0.6);
\node[font=\tiny,below] at (3.4,-0.6) {dwell time (ціль: скоротити)};
\draw (6.5,0.1)--(6.5,-0.8); \node[font=\tiny,rotate=25,anchor=west] at (6.5,0.2) {виявлення};
\end{tikzpicture}
\figcap{Рис.~\thechapter.1 --- Час присутності (dwell time) як проміжок від компрометації до виявлення}
\end{tikzfig}

Зв'язок збитків із часом присутності часто монотонно зростаючий: що довше супротивник діє, то більше вузлів компрометує, то більше даних викрадає чи руйнує. Спрощену залежність очікуваних збитків $D$ від dwell time $T$ можна подати як

\begin{equation}
D(T) = D_0\bigl(e^{\gamma T} - 1\bigr),
\label{eq:ch13-damage}
\end{equation}

де $D_0$ й $\gamma>0$ --- параметри, що відображають початковий темп і прискорення шкоди (бокове переміщення розширює доступ, збільшуючи й темп). Експоненційний характер~\eqref{eq:ch13-damage} пояснює, чому навіть помірне скорочення dwell time дає непропорційно велике зниження збитків, і робить раннє виявлення --- а отже полювання --- економічно виправданим попри його вартість.

Автоматичне виявлення скорочує dwell time для \emph{відомих} загроз, але APT навмисно їх уникає. Тому для скорочення dwell time проти невідомого потрібен проактивний пошук --- полювання. Емпірично організації з розвиненим полюванням виявляють APT значно раніше, ніж ті, що покладаються лише на реактивні сповіщення.

\section{Парадигма полювання за загрозами}\label{sec:ch13-hunting}

\term{Полювання за загрозами}~\cite{ch13:sqrrl} --- проактивний, ітеративний, керований людиною пошук ознак компрометації, що уникли автоматичного виявлення. Його відмінність від реагування на сповіщення: полювання починається не з тривоги, а з \emph{гіпотези} («якщо угруповання $X$ діє в нашій мережі, ми побачимо ознаку $Y$»), яку аналітик перевіряє в даних. Полювання виходить із припущення компрометації (assume breach): замість питати «чи спрацювала тривога?», аналітик питає «що ми могли пропустити?».

\begin{tikzfig}
\begin{tikzpicture}[font=\scriptsize,
  b/.style={draw,rounded corners=2pt,minimum height=5.5mm,text width=18mm,align=center,font=\tiny},
  >={Stealth[length=1.4mm]}]
\node[b] (h) at (0,0) {гіпотеза (розвідка, TTP)};
\node[b] (d) at (2.4,0) {пошук у даних};
\node[b] (f) at (4.8,0) {знахідка?};
\node[b] (r) at (7.0,0.6) {реагування};
\node[b] (a) at (7.0,-0.6) {нове правило виявлення};
\draw[->] (h)--(d); \draw[->] (d)--(f); \draw[->] (f)--(r); \draw[->] (f)--(a);
\draw[->] (a) to[out=-150,in=-30] (h);
\end{tikzpicture}
\figcap{Рис.~\thechapter.2 --- Цикл полювання за загрозами: від гіпотези до реагування й нового автоматизованого виявлення}
\end{tikzfig}

Цінний продукт полювання --- не лише знайдені загрози, а й \emph{автоматизація}: успішна гіпотеза стає новим правилом виявлення (розділ~\ref{chap:ch11}), тож одноразовий ручний пошук перетворюється на постійне автоматичне виявлення. Так полювання й інженерія виявлення утворюють контур: полювання знаходить нове, виявлення автоматизує знайдене, звільняючи мисливців для пошуку наступного невідомого.

Полювання принципово спирається на \emph{людину}, і це не тимчасове обмеження, а сутнісна риса. Автоматика чудова в розпізнаванні \emph{відомих} патернів, але погано справляється з нечіткими, контекстуальними, уперше баченими ознаками, де потрібні інтуїція, доменне знання й творче формулювання гіпотез. Аналітик-мисливець привносить саме це: розуміння, що \emph{мало б} статися за даної гіпотези, навіть якщо точного правила ще не існує. Тому полювання не конкурує з автоматизацією, а працює на її межі --- там, де сигнали занадто слабкі чи незвичні для наявних правил, --- і поступово розширює автоматизовану зону, перетворюючи вчорашні відкриття на сьогоднішні правила. Ця взаємодія людини й машини --- визначальна риса зрілого кіберзахисту проти мислячого супротивника.

\section{Методологія й гіпотези}\label{sec:ch13-methodology}

Гіпотези полювання походять із трьох джерел: \emph{розвідки загроз} (розділ~\ref{chap:ch12}) --- TTP відомих угруповань, релевантних моделі загроз; \emph{аномалій} --- статистичних відхилень у телеметрії; \emph{обізнаності про середовище} --- знання, що є нормою, і що --- ні. Методично полювання спирається на багату телеметрію (розділ~\ref{chap:ch10}) й аналітику; репозиторії аналітик на кшталт MITRE CAR~\cite{mitre-car} надають готові запити для виявлення технік. Модель «хто, що, де, коли, чому, як»~\cite{ch13:ocm} структурує планування полювання.

Практично гіпотези полювання доцільно формулювати так, щоб вони були \emph{спростовними} й \emph{перевірними} наявними даними. Добра гіпотеза вказує конкретну очікувану ознаку та джерело, де її шукати: не розпливчасте «чи є в нас APT?», а «якщо застосовано техніку $T$, то в телеметрії створення процесів за останні 30~діб ми побачимо батьківсько-дочірній ланцюг виду $P$». Така постановка дозволяє або знайти підтвердження, або впевнено його відкинути, або виявити, що потрібних даних немає (прогалина видимості). Останній результат особливо цінний: він перетворює полювання на інструмент оцінювання якості телеметрії, спрямовуючи інвестиції в ті джерела подій, яких бракує для перевірки пріоритетних гіпотез. Відсутність даних --- тиха, але небезпечна форма сліпоти, яку реактивний захист не помічає, бо не ставить питань, на які немає відповіді.

Окремий методологічний принцип --- зосередження на техніках, стійких у піраміді болю (розділ~\ref{chap:ch11}). Полювання за ефемерними індикаторами (адресами, гешами) малоцінне: їх і так покриває автоматика, а зловмисник легко змінює. Натомість полювання за поведінковими TTP --- закономірностями, які угруповання змушене повторювати через обмеженість способів досягнення мети, --- дає стійкі результати й породжує довговічні правила. Тому зрілі програми пріоритезують гіпотези навколо технік, від яких супротивникові найважче відмовитися: механізмів закріплення, бокового переміщення, вивантаження даних.

Ефективність полювання зручно оцінити через приріст виявлюваності. Якщо автоматичне виявлення покриває техніки з повнотою $r_{\text{auto}}$, а полювання додає виявлення технік, що уникли автоматики, з повнотою $r_{\text{hunt}}$ на решті, то сукупна повнота

\begin{equation}
r_{\Sigma} = r_{\text{auto}} + (1 - r_{\text{auto}})\,r_{\text{hunt}},
\label{eq:ch13-hunt}
\end{equation}

перевищує будь-яку окремо. Формула~\eqref{eq:ch13-hunt} формалізує взаємодоповнення: полювання цінне саме там, де автоматика безсила ($r_{\text{auto}}$ мала для нових TTP), тому інвестиція в полювання дає найбільший приріст проти витончених, невідомих загроз --- тобто проти APT.

\section{Бокове переміщення та living-off-the-land як об'єкти полювання}\label{sec:ch13-lateral}

Найхарактерніші й найважчі для автоматичного виявлення дії APT --- \term{бокове переміщення} (поширення мережею від початкового вузла до цілей) та зловживання легітимними інструментами (LotL, розділ~\ref{chap:ch11}; в Unix-середовищах --- GTFOBins~\cite{ch13:gtfobins}). Бокове переміщення використовує легітимні механізми (віддалене виконання, викрадені облікові дані, адміністративні протоколи), тому зливається з нормальною активністю й погано ловиться сигнатурами. Його виявлення~\cite{ch13:lateral} спирається на аномалії: нетипові зв'язки між вузлами, незвичні часи, аномальне використання облікових записів.

Полювання за цими діями --- показовий приклад переваги гіпотезо-керованого підходу: замість чекати сигнатури аналітик висуває гіпотезу («скомпрометований обліковий запис використовується для доступу до нетипових вузлів») і шукає її підтвердження в графі автентифікацій і з'єднань. Ключ --- модель норми: що є типовим для цього користувача, вузла, часу. Відхилення від норми, контекстуалізоване розвідкою й знанням середовища, і є здобиччю полювання. Українські практики розглядають полювання як складову стратегії кіберзахисту~\cite{ch13:hunt-ua} та методи виявлення цілеспрямованих атак на критичну інфраструктуру~\cite{ch13:apt-crit-ua}.

\section{Зрілість і метрики полювання}\label{sec:ch13-maturity}

Зрілість полювання описують моделлю~\cite{ch13:hunting-maturity}: від нульового рівня (покладання лише на автоматичні сповіщення) через ситуативне й кероване розвідкою полювання до найвищого --- автоматизованого, де успішні гіпотези систематично перетворюються на виявлення. Прогрес залежить від якості телеметрії, навичок аналітиків і процесів. Рівні зрілості узагальнено в табл.~\thechapter.1.

\tabhead{Таблиця~\thechapter.1 --- Модель зрілості полювання за загрозами}

\begin{center}\footnotesize
\begin{tabularx}{\linewidth}{L{16mm} Y L{20mm}}
\toprule
Рівень & Характеристика & Основа \\
\midrule
Початковий & лише автоматичні сповіщення & немає полювання \\
Мінімальний & включення індикаторів розвідки & зовнішня розвідка \\
Процедурний & регулярні процедури полювання & телеметрія, аналітики \\
Інноваційний & власні аналітики й гіпотези & розвинені навички \\
Провідний & автоматизація успішних гіпотез & контур із виявленням \\
\bottomrule
\end{tabularx}
\end{center}


Метрики полювання відрізняються від метрик виявлення: важлива не лише кількість знайдених загроз (успішне полювання може нічого не знайти, але підтвердити відсутність загрози й виявити прогалини видимості), а й \emph{приріст покриття} (скільки нових правил породжено), \emph{виявлені сліпі зони} (прогалини телеметрії), \emph{скорочення dwell time}. Парадокс полювання: «порожнє» полювання не є провалом --- воно або підтверджує безпеку, або виявляє, що бракує даних для перевірки гіпотези, що саме є цінним результатом. Тому полювання оцінюють за систематичністю й покриттям гіпотез, а не лише за «уловом».

\section{Кейс: полювання за прихованим угрупованням}\label{sec:ch13-case}

Розглянемо полювання за ознаками APT, що уникло автоматичного виявлення (захисна задача).

\emph{Крок 1. Гіпотеза.} З розвідки загроз (розділ~\ref{chap:ch12}) відомо, що релевантне угруповання застосовує бокове переміщення через викрадені адміністративні облікові дані. Формулюємо гіпотезу: «якщо угруповання діє, ми побачимо нетипові віддалені входи адміністративним обліковим записом».

\emph{Крок 2. Дані.} Збираємо телеметрію автентифікацій, створення процесів і мережевих з'єднань (розділ~\ref{chap:ch10}) за достатній період.

\emph{Крок 3. Пошук аномалій.} Будуємо модель норми (типові вузли, часи, облікові записи) і шукаємо відхилення: вхід адміністратора на вузли, куди він ніколи не заходив, у неробочий час, з подальшим запуском інтерпретатора команд.

\emph{Крок 4. Підтвердження й розширення.} Знайшовши підозрілий патерн, розширюємо розслідування діамантовою моделлю (розділ~\ref{chap:ch12}): від скомпрометованого вузла --- до інфраструктури, інших уражених вузлів, часу початкової компрометації (оцінка dwell time).

\emph{Крок 5. Реагування й автоматизація.} Передаємо підтверджену загрозу на реагування (розділ~\ref{chap:ch16}), а гіпотезу перетворюємо на правило виявлення (розділ~\ref{chap:ch11}), щоб надалі ловити цей патерн автоматично~\eqref{eq:ch13-hunt}.

Кейс ілюструє силу проактивного, гіпотезо-керованого пошуку проти супротивника, який навмисно уникає автоматичного виявлення: полювання скорочує dwell time, виявляє приховане й перетворює знахідки на постійне автоматизоване виявлення.

\section*{Висновки до розділу}
\addcontentsline{toc}{section}{Висновки до розділу}

Цілеспрямована стала загроза (APT) --- добре ресурсоване, наполегливе, адаптивне угруповання, що досягає конкретних цілей попри захист, уникаючи відомих сигнатур і зловживаючи легітимними інструментами. Проти APT реактивне виявлення недостатнє, бо такий супротивник діє приховано й може роками лишатися непоміченим, збільшуючи час присутності (dwell time), з яким зростають збитки. Відповіддю є полювання за загрозами --- проактивний, ітеративний, керований гіпотезами пошук прихованого, що виходить із припущення компрометації. Гіпотези походять із розвідки загроз, аномалій і знання середовища; цінний продукт полювання --- не лише знахідки, а й автоматизація успішних гіпотез у правила виявлення, що замикає контур із інженерією виявлення~\eqref{eq:ch13-hunt}. Найхарактерніші об'єкти полювання --- бокове переміщення й LotL, що зливаються з нормою й виявляються за аномаліями відносно моделі норми. Зрілість полювання зростає від реактивності до автоматизації, а його метрики оцінюють систематичність і покриття гіпотез, а не лише «улов», бо «порожнє» полювання теж цінне. Головний висновок: проти витончених невідомих загроз проактивне полювання є необхідним доповненням автоматичного виявлення, що найбільше скорочує dwell time саме там, де автоматика безсила.

\selfcheck
\begin{enumerate}
\item Якими трьома ознаками характеризують APT? Чим APT відрізняється від масового ЗПК?
\item Чому типовий периметровий захист і сигнатурне виявлення систематично пропускають APT?
\item Що таке час присутності (dwell time)? Чому його мінімізація є стратегічною метою?
\item Чому для скорочення dwell time проти невідомих загроз потрібен проактивний пошук?
\item Що таке полювання за загрозами? Чим воно відрізняється від реагування на сповіщення?
\item Поясніть припущення компрометації (assume breach) як основу полювання.
\item Чому автоматизація успішних гіпотез є цінним продуктом полювання? Як утворюється контур із інженерією виявлення?
\item З яких джерел походять гіпотези полювання? Наведіть приклад гіпотези.
\item Поясніть формулу сукупної повноти~\eqref{eq:ch13-hunt}. Чому полювання дає найбільший приріст проти APT?
\item Що таке бокове переміщення й LotL? Чому вони важкі для автоматичного виявлення?
\item Як виявляють бокове переміщення? Яку роль відіграє модель норми?
\item Опишіть модель зрілості полювання. Чому «порожнє» полювання не є провалом?
\item Опишіть полювання за прихованим угрупованням (кейс підрозділу~\ref{sec:ch13-case}). Як оцінюють dwell time?
\item Дослідницьке: розробіть і перевірте набір гіпотез полювання за технікою бокового переміщення на реалістичній телеметрії, оцініть виявлені сліпі зони.
\item Дослідницьке: оцініть вплив розвиненості полювання на dwell time, порівнявши сценарії суто реактивного та проактивного захисту.
\end{enumerate}
